\chapter*{第 1 章 绪论}\addcontentsline{toc}{chapter}{第 1 章 绪论}\markboth{第 1 章 绪论}{}

\begin{tishi}
\textbf{章号与题源说明：}本册按\textbf{赵琴教材第 1--8 章}归章，本章为\textbf{赵琴第 1 章 绪论}。
孔珑原书的对应关系是：\textbf{孔珑第 1 章 绪论（原书无课后习题）＋孔珑第 2 章 流体及其物理性质}
$\to$ \textbf{赵琴第 1 章}。因此本章收录的题目全部来自\textbf{孔珑教材第 2 章"流体及其物理性质"
的课后习题}，保留孔珑原书题号 \textbf{2-1 至 2-17}（共 17 题，\textbf{一题不落全部收录}）。

\textbf{本册星级规则：}每题星级取该题\textbf{最主要考点}在《教材精讲》第 1 章中的星级
（密度、比容与重度＝3 星；压缩性与膨胀性＝3 星；黏性与牛顿内摩擦定律＝5 星；
表面张力＝2 星），以保证同一考点在《教材精讲》与本册中星数一致；
题中夹带的超过 818 考纲的经验公式（混合气体黏度、等熵体积模量、苏士兰式、恩氏黏度换算）
不改变星级，另在章末统一提示。

\textbf{做题要求：}计算题必须写出"公式—代入—单位—结果"四步；单位一律化为国际单位。
\end{tishi}

\begin{ti}
\sloppy
\emergencystretch=2em

\item \textbf{孔珑 2-1}\quad 已知某种物质的密度 $\rho=2.94\ \mathrm{g/cm^3}$，试求它的相对密度 $d$（取水的密度 $\rho_w=1000\ \mathrm{kg/m^3}$）。\xstarfull{3}\yiju{E4 2015 年试题填空第 1 题（同一"密度—重度—比容"换算考点）；E6 教材式 1.1--1.4；《教材精讲》第 1 章"密度、比容与重度"定为三星}\nandu{易}

\item \textbf{孔珑 2-2}\quad 已知某厂 1 号炉水平烟道中烟气各组分的体积分数为 $\alpha_{\mathrm{CO_2}}=13.5\%$、$\alpha_{\mathrm{SO_2}}=0.3\%$、$\alpha_{\mathrm{O_2}}=5.2\%$、$\alpha_{\mathrm{N_2}}=76\%$、$\alpha_{\mathrm{H_2O}}=5\%$，试求该烟气的密度（各组分在标准状态下的密度分别为 $1.967$、$2.927$、$1.429$、$1.251$、$0.804\ \mathrm{kg/m^3}$）。\xstarfull{3}\yiju{E4 2022 年 818 单选第 4 题（压缩机进出口气体密度计算）；E6 教材式 1.1、1.2 与混合气体密度求法；《教材精讲》第 1 章"密度、比容与重度"定为三星}\nandu{中}

\item \textbf{孔珑 2-3}\quad 在孔珑 2-2 题中，实测烟气温度 $t=170\ ^{\circ}\mathrm{C}$、静计示压强 $p_e=1432\ \mathrm{Pa}$、当地大气压强 $p_a=100858\ \mathrm{Pa}$。试求\textbf{工作状态下}烟气的密度和运动黏度。\xstarfull{3}\yiju{第一问（密度）依据 E4 2022 年 818 单选第 4 题与 E6 教材式 1.1、1.2、状态方程 $\rho=p/(RT)$，对应《教材精讲》第 1 章"密度、比容与重度"三星；第二问（混合气体运动黏度）所需的苏士兰式与 Wilke 混合式为超纲扩展公式，不计入星级，详见章末提示}\nandu{难}

\item \textbf{孔珑 2-4}\quad 当压强增量为 $50000\ \mathrm{Pa}$ 时，某种液体的密度增长 $0.02\%$，试求该液体的体积模量。\xstarfull{3}\yiju{E5 题库判断题"体积模量 $K$ 值越大，液体越容易压缩"（同一考点）；E6 教材式 1.5--1.7；《教材精讲》第 1 章"流体的压缩性与膨胀性"定为三星}\nandu{易}

\item \textbf{孔珑 2-5}\quad 已知气体的压强由 $p_1=1.01325\times10^{5}\ \mathrm{Pa}$ 增至 $p_2=3.923\times10^{5}\ \mathrm{Pa}$，试求该气体在\textbf{等温过程}和\textbf{等熵过程}中的体积模量（等熵指数取 $\kappa=1.4$）。\xstarfull{3}\yiju{E6 教材式 1.5--1.7 与教材 1.3.2 节"气体等温、等熵过程的体积模量"；E5 题库同类概念题；《教材精讲》第 1 章"流体的压缩性与膨胀性"定为三星}\nandu{中}
\tuyi{本题题面在原书 OCR 中只残存末尾"多少？"三字，题面中的已知压强 $p_1$、$p_2$ 由原书解答反推补全（数值与书末答案 $3.923\times10^{5}\ \mathrm{Pa}$、$5.492\times10^{5}\ \mathrm{Pa}$ 完全吻合），详见章末"题面 OCR 不清"说明。}

\item \textbf{孔珑 2-6}\quad 充满石油的油槽内的压强为 $4.9033\times10^{5}\ \mathrm{Pa}$，今由槽中排出石油 $40\ \mathrm{kg}$，使槽内压强降到 $9.8067\times10^{4}\ \mathrm{Pa}$。设石油的密度 $\rho=880\ \mathrm{kg/m^3}$、体积模量 $K=1.32\times10^{9}\ \mathrm{Pa}$，试求油槽的体积。\xstarfull{3}\yiju{E6 教材式 1.5--1.7（由压缩性定义反求体积）；《教材精讲》第 1 章"流体的压缩性与膨胀性"定为三星；E4 历年真题中未见独立成题}\nandu{中}

\item \textbf{孔珑 2-7}\quad 流量为 $50\ \mathrm{m^3/h}$、温度为 $70\ ^{\circ}\mathrm{C}$ 的水流入热水锅炉，经加热后水温升到 $90\ ^{\circ}\mathrm{C}$，已知水的体积膨胀系数 $\alpha_V=0.00064\ /^{\circ}\mathrm{C}$，问从锅炉中每小时流出多少立方米的水？\xstarfull{3}\yiju{E6 教材式 1.8（体积膨胀系数的定义）；E5 题库"流体的物理特性有密度、比容和体积膨胀系数"；《教材精讲》第 1 章"流体的压缩性与膨胀性"定为三星}\nandu{易}

\item \textbf{孔珑 2-8}\quad 压缩机压缩空气，绝对压强从 $9.8067\times10^{4}\ \mathrm{Pa}$ 升高到 $5.8840\times10^{5}\ \mathrm{Pa}$，温度从 $20\ ^{\circ}\mathrm{C}$ 升高到 $78\ ^{\circ}\mathrm{C}$，问空气体积减小了多少？\xstarfull{3}\yiju{E4 2022 年 818 单选第 4 题（空气进压缩机，求进出口密度与质量流量，同一考点）；E6 教材式 1.1、1.2 与状态方程；《教材精讲》第 1 章"密度、比容与重度"定为三星}\nandu{中}

\item \textbf{孔珑 2-9}\quad 动力黏度为 $2.9\times10^{-4}\ \mathrm{Pa\cdot s}$、密度为 $678\ \mathrm{kg/m^3}$ 的油，其运动黏度等于多少？\xstarfull{5}\yiju{E4 2015 年试题填空第 1 题（由密度与运动黏度求重度与动力黏度，同一 $\nu=\mu/\rho$ 换算）；E5 题库 17.pdf 第 8 题（$\nu=3\times10^{-6}\ \mathrm{m^2/s}$、$\rho=800\ \mathrm{kg/m^3}$ 求 $\mu$）；E6 教材式 1.11；《教材精讲》第 1 章"黏性"定为五星}\nandu{易}

\item \textbf{孔珑 2-10}\quad 设空气在 $0\ ^{\circ}\mathrm{C}$ 时的运动黏度 $\nu_0=13.2\times10^{-6}\ \mathrm{m^2/s}$、密度 $\rho_0=1.29\ \mathrm{kg/m^3}$，试求在 $150\ ^{\circ}\mathrm{C}$ 时空气的动力黏度（提示：查教材表 2-6 得苏士兰常数 $S=111\ \mathrm{K}$）。\xstarfull{5}\yiju{E4 2022 年 818 判断题第 6 题（气体与液体黏度随温度变化的趋势及原因）；E6 教材 2.3 节"黏度与温度的关系"；《教材精讲》第 1 章"黏性"定为五星（苏士兰式为超纲经验式，见章末提示）}\nandu{中}

\item \textbf{孔珑 2-11}\quad 借恩氏黏度计测得石油的黏度为 $8.5\ ^{\circ}\mathrm{E}$，如石油的密度 $\rho=850\ \mathrm{kg/m^3}$，求石油的动力黏度。\xstarfull{2}\yiju{E4 历年真题与 E5 题库中均未见恩氏黏度换算题；E6 教材 2.3 节黏度测量属扩展内容，《教材精讲》第 1 章未单列该考点；按本册星级表"真题未出现且非考纲核心"定为二星}\nandu{易}

\item \textbf{孔珑 2-12}\quad 一平板距离另一固定平板 $0.5\ \mathrm{mm}$，两板间充满液体，上板在每平方米有 $2\ \mathrm{N}$ 的力作用下以 $0.25\ \mathrm{m/s}$ 的速度移动，求该液体的黏度。
\tuyi{两平行平板，下板固定，两板间距 $h=0.5\ \mathrm{mm}$；上板在单位面积受力 $\tau=2\ \mathrm{N/m^2}$ 的作用下以 $U=0.25\ \mathrm{m/s}$ 匀速平动，板间液体速度按线性分布。}
\xstarfull{5}\yiju{E4 2026 年回忆版计算题第 1 题（牛顿流体、切应力问题）；E4 2022 年 818 判断题第 2 题（同轴圆柱间隙内筒表面切应力 $\tau=\mu U/h$）；E6 教材式 1.9、1.10；《教材精讲》第 1 章"黏性"定为五星}\nandu{易}

\item \textbf{孔珑 2-13}\quad 已知动力滑动轴承的轴直径 $d=0.2\ \mathrm{m}$，转速 $n=2830\ \mathrm{r/min}$，轴承内径 $D=0.2016\ \mathrm{m}$，宽度 $l=0.3\ \mathrm{m}$，润滑油的动力黏度 $\mu=0.245\ \mathrm{Pa\cdot s}$，试求克服摩擦阻力所消耗的功率。
\tuyi{滑动轴承的轴与轴瓦之间充满润滑油，轴径 $d=0.2\ \mathrm{m}$、轴瓦内径 $D=0.2016\ \mathrm{m}$，径向间隙 $h=(D-d)/2=0.0008\ \mathrm{m}$；轴以 $n=2830\ \mathrm{r/min}$ 旋转，受剪面积 $A=\pi d l$。}
\xstarfull{5}\yiju{E4 2026 年回忆版计算题第 1 题（牛顿流体黏性阻力）；E4 2022 年 818 判断题第 2 题（间隙内 $\tau=\mu U/h$）；E6 教材式 1.9、1.10；《教材精讲》第 1 章"黏性"定为五星}\nandu{中}

\item \textbf{孔珑 2-14}\quad 某转子以 $600\ \mathrm{r/min}$ 旋转时，它的角减速度为 $0.02\ \mathrm{rad/s^2}$。已知轴套的长度为 $5\ \mathrm{cm}$，轴的直径为 $2\ \mathrm{cm}$，以及它们之间的间隙为 $0.05\ \mathrm{mm}$，试求流体的黏度（转子的转动惯量 $J=50\times0.3^{2}=4.5\ \mathrm{kg\cdot m^2}$）。
\tuyi{轴套静止，轴在其中转动；轴直径 $d=2\ \mathrm{cm}$，轴套长度 $l=5\ \mathrm{cm}$，轴与轴套之间的径向间隙 $h=0.05\ \mathrm{mm}$，间隙内充满待测流体；轴以 $n=600\ \mathrm{r/min}$ 转动，因流体黏性阻力产生角减速度 $\varepsilon=0.02\ \mathrm{rad/s^2}$。}
\xstarfull{5}\yiju{E4 2022 年 818 判断题第 2 题（同轴圆柱间隙内筒表面切应力 $\tau=\mu U/h$）；E6 教材式 1.9、1.10；《教材精讲》第 1 章"黏性"定为五星}\nandu{中}

\item \textbf{孔珑 2-15}\quad 直径为 $5.00\ \mathrm{cm}$ 的活塞在直径为 $5.01\ \mathrm{cm}$ 的缸体内运动。当润滑油的温度由 $0\ ^{\circ}\mathrm{C}$ 升高到 $120\ ^{\circ}\mathrm{C}$ 时，求推动活塞所需的力减少的百分数（用教材图 2-5 中相对密度 $d=0.855$ 的原油的黏度进行计算）。
\tuyi{活塞与缸体之间的环形间隙内充满润滑油，活塞直径 $5.00\ \mathrm{cm}$、缸体内径 $5.01\ \mathrm{cm}$，径向间隙 $h=0.005\ \mathrm{cm}$；活塞在缸体内沿轴向运动，润滑油温度由 $0\ ^{\circ}\mathrm{C}$ 升到 $120\ ^{\circ}\mathrm{C}$，动力黏度由 $\mu_1$ 降到 $\mu_2$。}
\xstarfull{5}\yiju{E4 2022 年 818 判断题第 6 题（黏度主要受温度影响）；E6 教材图 2-5 与式 1.9、1.10；《教材精讲》第 1 章"黏性"定为五星}\nandu{中}

\item \textbf{孔珑 2-16}\quad 内径为 $10\ \mathrm{mm}$ 的开口玻璃管插入温度为 $20\ ^{\circ}\mathrm{C}$ 的水中，已知水与玻璃的接触角 $\theta=10^{\circ}$，试求水在管中上升的高度。
\tuyi{竖直开口玻璃管插入水中，管内径 $d=10\ \mathrm{mm}$；水能润湿玻璃，管内液面呈凹形并高出管外自由液面，液面切线与管壁的接触角 $\theta=10^{\circ}$。}
\xstarfull{2}\yiju{E4 历年真题与 E5 题库中未见毛细现象成题；E6 教材 1.3.4 节，教材原文明言"由于表面张力很小，一般对液体的宏观运动不起作用，可以忽略"；《教材精讲》第 1 章"液体的表面张力"定为二星}\nandu{易}

\item \textbf{孔珑 2-17}\quad 内径 $8\ \mathrm{mm}$ 的开口玻璃管插入 $20\ ^{\circ}\mathrm{C}$ 的水银中，已知水银与玻璃的接触角约为 $140^{\circ}$，试求水银在管中下降的高度（查教材表 2-9 得 $20\ ^{\circ}\mathrm{C}$ 水银与空气接触的表面张力系数 $\sigma=0.5137\ \mathrm{N/m}$，水银密度 $\rho=13550\ \mathrm{kg/m^3}$）。
\tuyi{竖直开口玻璃管插入水银中，管内径 $d=8\ \mathrm{mm}$；水银不能润湿玻璃，管内液面呈凸形并低于管外自由液面（下降），接触角 $\theta=140^{\circ}>90^{\circ}$。}
\xstarfull{2}\yiju{E4 历年真题与 E5 题库中未见毛细现象成题；E6 教材 1.3.4 节与表 2-9；《教材精讲》第 1 章"液体的表面张力"定为二星}\nandu{易}

\end{ti}

\begin{tishi}
\textbf{本章略去的题号及原因：无。}
孔珑第 2 章课后习题共 17 题（2-1 至 2-17），全部属于"流体的主要物理性质"，
对应赵琴教材第 1 章（绪论），\textbf{全部符合 818 考纲，故一题不落全部收录}。
本章没有出现气体动力学（孔珑第 7、10 章）、流体机械（泵与风机）、明渠堰流等超纲题，
不存在因超纲而略去的题号。
\end{tishi}

\begin{tishi}
\textbf{OCR 校订清单（明显不合理的数值按物理量纲订正）：}
（1）孔珑 2-4 原书答案 OCR 作 "$2.5\times10^{2}\ \mathrm{Pa}$"，按
$K=\Delta p/(\Delta\rho/\rho)$ 应为 $2.5\times10^{8}\ \mathrm{Pa}$；
（2）孔珑 2-6 的油槽压强 OCR 作 "$4.9033\times10^{6}\ \mathrm{Pa}$"，按书末答案
$V=153\ \mathrm{m^3}$ 反推应为 $4.9033\times10^{5}\ \mathrm{Pa}$；
（3）多处密度单位 OCR 作 $\mathrm{kg/m^2}$、体积单位作 $\mathrm{m^2}$，
一律订正为 $\mathrm{kg/m^3}$、$\mathrm{m^3}$（涉及 2-2、2-3、2-6、2-7、2-9、2-17）；
（4）孔珑 2-10 题给"密度 $\rho_0=2.94\ \mathrm{g/cm^2}$"OCR 有误，原书解答中实际取
$\rho_0=1.29\ \mathrm{kg/m^3}$（$0\ ^{\circ}\mathrm{C}$ 空气密度），题面按此订正；
（5）孔珑 2-13 的功率 OCR 作 "$50.7\times10^{*}\ \mathrm{W}$"，按计算趋势订正为
$50.7\times10^{3}\ \mathrm{W}$；
（6）孔珑 2-15 图中读数 OCR 把两个指数都印成 $10^{-2}$，因题目问"力\textbf{减少}"必有
$\mu_2<\mu_1$，故 $120\ ^{\circ}\mathrm{C}$ 时的黏度订正为 $2.2\times10^{-3}\ \mathrm{Pa\cdot s}$；
（7）孔珑 2-7 的体胀系数 OCR 作 "$0.000641/^{\circ}\mathrm{C}$"，应读作
$\alpha_V=0.00064\ /^{\circ}\mathrm{C}$。
\end{tishi}

\begin{tishi}
\textbf{题面 OCR 不清、按标准解法处理的三处说明：}
（1）\textbf{孔珑 2-5}：题面 OCR 只残存末尾"多少？"三字，本册按原书解答反推补全为
"已知 $p_1=1.01325\times10^{5}\ \mathrm{Pa}$、$p_2=3.923\times10^{5}\ \mathrm{Pa}$，
求等温与等熵过程的体积模量"，补全后的数值与原书答案
$3.923\times10^{5}\ \mathrm{Pa}$、$5.492\times10^{5}\ \mathrm{Pa}$ 完全吻合；
（2）\textbf{孔珑 2-14}：题面开头 OCR 缺失，本册按原书解答中的
$J\varepsilon=50\times0.3^{2}\times0.02=0.09\ \mathrm{N\cdot m}$ 反推，
补入"转子的转动惯量 $J=50\times0.3^{2}=4.5\ \mathrm{kg\cdot m^2}$"；
（3）\textbf{孔珑 2-16}：原书解答在 OCR 中未保留，本册按标准解法处理。
\end{tishi}

\begin{tishi}
\textbf{超出 818 考纲的扩展内容提示（题目仍收录，但相应的经验公式可不记）：}
（1）\textbf{孔珑 2-3} 第二问的\textbf{混合气体动力黏度}：既要用苏士兰式
$\mu=\mu_0\dfrac{273+S}{T+S}\left(\dfrac{T}{273}\right)^{3/2}$，又要用 Wilke 混合式
$\mu=\dfrac{\sum\alpha_i M_i^{1/2}\mu_i}{\sum\alpha_i M_i^{1/2}}$，两者都超出考纲；
只需掌握第一问"由状态方程求工作状态密度"和"由 $\nu=\mu/\rho$ 求运动黏度"的思路；
（2）\textbf{孔珑 2-5} 的\textbf{等熵}体积模量 $K=\kappa p$（$\kappa=1.4$）超出考纲，
等温过程的 $K=p$ 属教材正文内容；
（3）\textbf{孔珑 2-10} 的\textbf{苏士兰式}超出考纲，但"气体黏度随温度升高而增大"是必背结论
（见《教材精讲》第 1 章"黏性"考点，2022 年 818 判断题第 6 题即考此结论）；
（4）\textbf{孔珑 2-11} 的\textbf{恩氏黏度换算}式
$\nu=0.0731\ ^{\circ}\mathrm{E}-0.0631/^{\circ}\mathrm{E}$ 超出考纲，仅作了解。
\end{tishi}
